\documentclass[11pt,a4paper]{article}
\input{../template/utm-lab}

\newcommand{\labnumber}{10}
\newcommand{\labtitle}{Add Your Own System Call: trace}
\newcommand{\studentname}{Daniel Nicorici}
\newcommand{\studentgroup}{FAF-243}

\begin{document}

\utmtitlepage
\tableofcontents
\newpage

\plainsection{Purpose of the work}

In this lab the kernel itself is changed for the first time. The purpose is to add a new system
call to xv6, \code{trace}, which makes a process print every system call it makes, the same
thing \code{strace} did on Linux in Lab 9, and to follow one number, 23, from a user program
down to a function in the kernel.

The work was done in my \code{\textasciitilde/xv6-riscv} tree (\code{mit-pdos/xv6-riscv}, commit
\code{06aad25}), cross-compiled with \code{riscv64-linux-gnu-gcc} 16.2.1 and run on QEMU 11.1.1.
In this tree number 22 is already \code{sync} and the kernel prints with \code{printk}, so the
new call uses 23 and \code{printk}.

\section{Part 1. The number and the user side}

Five small edits give the call a number and teach the user side to ask for it.

\begin{lstlisting}[caption={The number, the declaration, the stub and the UPROGS line}]
$ git diff kernel/syscall.h user/user.h user/usys.pl
--- a/kernel/syscall.h
+++ b/kernel/syscall.h
@@ -21,3 +21,4 @@
 #define SYS_mkdir  20
 #define SYS_close  21
 #define SYS_sync   22
+#define SYS_trace  23
--- a/user/user.h
+++ b/user/user.h
@@ -25,6 +25,7 @@ char *sys_sbrk(int, int);
 int pause(int);
 int uptime(void);
 int sync(void);
+int trace(int);
 
 // ulib.c
 int stat(const char *, struct stat *);
--- a/user/usys.pl
+++ b/user/usys.pl
@@ -43,3 +43,4 @@ entry("sbrk");
 entry("pause");
 entry("uptime");
 entry("sync");
+entry("trace");
$ grep -n "_trace" Makefile
149:	$U/_trace\
\end{lstlisting}

At build time \code{usys.pl} writes \code{user/usys.S}, where \code{trace} becomes three
instructions: \code{li a7, SYS\_trace}, \code{ecall}, \code{ret}.

\begin{lstlisting}[style=ccode,caption={user/trace.c}]
#include "kernel/types.h"
#include "kernel/stat.h"
#include "user/user.h"

int
main(int argc, char *argv[])
{
  if (argc < 3) {
    fprintf(2, "usage: trace mask command [args...]\n");
    exit(1);
  }
  if (trace(atoi(argv[1])) < 0) {
    fprintf(2, "trace: failed\n");
    exit(1);
  }
  exec(argv[2], argv + 2);
  fprintf(2, "trace: exec %s failed\n", argv[2]);
  exit(1);
}
\end{lstlisting}

The program sets the mask on itself and then becomes the command with \code{exec}; the mask
stays, because \code{exec} keeps the same process.

\begin{lstlisting}[caption={Check 1: the number reaches the kernel}]
$ trace 32 grep hello README
3 trace: unknown sys call 23
trace: failed
\end{lstlisting}

The kernel has no function for 23 yet, so it complains, but the message proves that the program
ran \code{ecall} with \code{a7 = 23} and that the kernel received it.

\section{Part 2. The kernel side}

\begin{lstlisting}[style=ccode,firstnumber=103,caption={kernel/proc.h: a place for the mask in struct proc}]
  char name[16];               // Process name (debugging)
  int trace_mask;
\end{lstlisting}

\begin{lstlisting}[style=ccode,firstnumber=114,caption={kernel/sysproc.c: the handler}]
uint64
sys_trace(void)
{
  int mask;
  argint(0, &mask);
  myproc()->trace_mask = mask;
  return 0;
}
\end{lstlisting}

\begin{lstlisting}[style=ccode,caption={kernel/syscall.c: the table entry, the names and the print}]
extern uint64 sys_sync(void);
extern uint64 sys_trace(void);
...
  [SYS_sync]    = sys_sync,
  [SYS_trace]   = sys_trace,
...
static char *syscallnames[] = {
  [SYS_fork] "fork",   [SYS_exit] "exit",     [SYS_wait] "wait",
  [SYS_pipe] "pipe",   [SYS_read] "read",     [SYS_kill] "kill",
  [SYS_exec] "exec",   [SYS_fstat] "fstat",   [SYS_chdir] "chdir",
  [SYS_dup] "dup",     [SYS_getpid] "getpid", [SYS_sbrk] "sbrk",
  [SYS_pause] "pause", [SYS_uptime] "uptime", [SYS_open] "open",
  [SYS_write] "write", [SYS_mknod] "mknod",   [SYS_unlink] "unlink",
  [SYS_link] "link",   [SYS_mkdir] "mkdir",   [SYS_close] "close",
  [SYS_sync] "sync",   [SYS_trace] "trace",
};
...
    p->trapframe->a0 = syscalls[num]();
    if ((p->trace_mask >> num) & 1)
      printk("%d: syscall %s -> %ld\n",
             p->pid, syscallnames[num], p->trapframe->a0);
\end{lstlisting}

\begin{lstlisting}[style=ccode,firstnumber=290,caption={kernel/proc.c: children inherit the mask in fork()}]
  safestrcpy(np->name, p->name, sizeof(p->name));
  np->trace_mask = p->trace_mask;
\end{lstlisting}

The handler reads argument 0 (register \code{a0}) with \code{argint} and stores it on the
current process. The print is placed after the call, so the return value is already in
\code{a0} of the trapframe.

\begin{lstlisting}[caption={Check 2: trace works}]
$ trace 32 grep hello README
3: syscall read -> 1023
3: syscall read -> 965
3: syscall read -> 453
3: syscall read -> 0
\end{lstlisting}

The numbers are exactly the expected ones: \(1023 + 965 + 453 = 2441\), the size of
\code{README}, read in three pieces, then 0 for the end of the file.

\section{Part 3. Play with the mask}

The mask is a number read in binary: bit \(N\) on means ``print call number \(N\)''.

\begin{lstlisting}[caption={All the calls (all bits on)}]
$ trace 2147483647 grep hello README
4: syscall trace -> 0
4: syscall exec -> 3
4: syscall open -> 3
4: syscall read -> 1023
4: syscall read -> 965
4: syscall read -> 453
4: syscall read -> 0
4: syscall close -> 0
\end{lstlisting}

\begin{lstlisting}[caption={Only fork (mask 2), on a test that fills the process table}]
$ trace 2 usertests forkforkfork
usertests starting
test forkforkfork: 4: syscall fork -> 5
5: syscall fork -> 6
6: syscall fork -> 7
6: syscall fork -> 8
6: syscall fork -> 9
6: syscall fork -> 10
[... 55 more fork lines from pids 6, 7 and 8 ...]
7: syscall fork -> -1
6: syscall fork -> -1
8: syscall fork -> -1
OK
ALL TESTS PASSED
\end{lstlisting}

\begin{lstlisting}[caption={Only write (mask 65536)}]
$ trace 65536 echo hi
hi5: syscall write -> 2

5: syscall write -> 1
\end{lstlisting}

\begin{lstlisting}[caption={Only open and close (mask 2129920)}]
$ trace 2129920 grep hello README
3: syscall open -> 3
3: syscall close -> 0
\end{lstlisting}

\textbf{Observations:}
\begin{enumerate}
  \item \textbf{Why only \code{read} with mask 32.} \(32 = 2^5\) has only bit 5 on, and 5 is
  \code{SYS\_read}. \code{open} is call 15 and \code{close} is call 21; their bits are 0, so the
  test \code{(p->trace\_mask >> num) \& 1} is false for them and nothing is printed.
  \item \textbf{Why several pids print in \code{forkforkfork}.} Because children inherit the
  mask: the line \code{np->trace\_mask = p->trace\_mask;} in \code{fork()} copies it from the
  parent to the child. In my run pids 4 and 5 printed one line each and pids 6, 7 and 8 printed
  the rest, until \code{fork} returned \(-1\) when the process table was full.
  \item \textbf{Why the trace line is between ``hi'' and the new line.} \code{echo} calls
  \code{write} twice: once for the word \code{hi} (return value 2) and once for the newline
  (return value 1). The kernel prints the trace line right after each call returns, so the first
  trace line comes after \code{hi} and before the newline of the second \code{write}.
  \item \textbf{The mask for \code{open} and \code{close} only.} \code{open} is 15 and
  \code{close} is 21, so the mask is \(2^{15} + 2^{21} = 32768 + 2097152 = 2129920\). Figure 3.4
  confirms it: only \code{open -> 3} and \code{close -> 0} are printed.
\end{enumerate}

The optional \code{sysinfo} call was not done.

\plainsection{Conclusions}

Adding \code{trace} showed the whole path of a system call, because every station on that path
had to be edited by hand. The call needs a number in \code{kernel/syscall.h}, a declaration in
\code{user/user.h} and a stub in \code{usys.pl} that puts the number in register \code{a7} and
executes \code{ecall}. On the kernel side the number is only an index into the table of
\code{syscall.c}, and the handler reads its argument from the saved registers of the process.
The first check made this concrete: with only the user side in place, the kernel already
received 23 and answered that it did not know it.

The mask also showed where per-process state lives: in \code{struct proc}. It survives
\code{exec}, because \code{exec} changes the program and not the process, and it reaches the
children only because \code{fork} copies it explicitly. With about thirty lines of code, xv6
now has a small \code{strace} of its own.

\end{document}
